\Def Операция $f$ ассоциативна, если $f (a, f (b, c)) = f (f (a, b), c)$.

\Def Операция $f$ коммутативна, если $f (a, b) = f (b, a)$.

Сначала реализуем стек с поддержкой произвольной ассоциативной операции. Для этого стек должен хранить пары. Для наглядности рассмотрим пример абстрактного стека, который хранит нужные пары вида $\{a, b\}$, где $a$ - непосредственно сам элемент, $b$ - значение функция от $a$ и элементов "ниже"\ $a$ (Индекс обозначает порядок добавления в стек):

\begin{table}[h]
\centering
\begin{tabular}{|l|l|}
\hline
$el_4$ & $f(el_4, f(el_3, f(el_2, el_1)$  \\
\hline
$el_3$ & $f(el_3, f(el_2, el_1)$           \\
\hline
$el_2$ & $f(el_2, el_1)$                    \\
\hline
$el_1$ & $el_1$\\
\hline
\end{tabular}
\end{table}

Реализация поддерживания подобной структуры при классическом добавлении/удалении элемента из стека очевидна.
\vspace{1cm}


Перейдём от стека к очереди.

Для этого будем поддерживать 2 стека с подобной структурой: стека на вход stack\_in, стек на выход stack\_out. В
первый будем добавлять, а из второго — извлекать.
Если при извлечении stack\_out пуст, то перекидываем все элементы
из stack\_in в stack\_out. Разумеется, мы поддерживаем вышеуказанную структуру в обоих стеках после каждого их удалений/добавлений.\\
Чтобы узнать $f$ от всех элементов очереди, достаточно получить значение $f(b_1, b_2)$, где $b_1, b_2$ - 2-ые элементы "верхних"\ пар стеков stack\_in и stack\_out.

\pagebreak